\section{Quantitative names, moduli, guards, and residual producers}
\label{sec:quantitative}

Effective analysis and continuous mathematics can be handled in \(\iPRA\) only
through finite quantitative statements.  DARC does not introduce completed
real or complex objects as a new domain.  It compiles statements about
primitive-recursive approximation programs, explicit error bounds, and
requested finite precisions.

\subsection{Finite-precision contracts}

A typical approximation node for an object \(x\) supplies a rational or dyadic
code \(a_x(p,\vec u)\) and a correctness formula
\[
  \operatorname{Approx}_x(p,\vec u,a_x(p,\vec u)).
\]
The formula may mean an interval enclosure, a scaled integer inequality, or a
metric bound encoded without division.  The exact encoding is theorem-specific,
but it must be decidable or reduced to reflected arithmetic atoms when used in
an induction invariant.

The public theorem asks for a fixed or variable numerical precision \(p\).
For a fixed artifact instance, the demand coordinate is capped at the maximum
precision reachable from that query.  Transfer rules propagate required input
precisions and guards backward.

\subsection{Addition and subtraction}

If \(|x-a|\le 2^{-(p+1)}\) and
\(|y-b|\le 2^{-(p+1)}\), then
\[
  |(x+y)-(a+b)|\le 2^{-p}.
\]
Thus a safe transfer for addition is
\[
  \Prec(x+y)\ge p
  \Longrightarrow
  \Prec(x)\ge p+1,\quad\Prec(y)\ge p+1.
\]
The same applies to subtraction.  In integer-scaled encodings, the proof is a
finite inequality after clearing powers of two.  A sharper local theorem may
reduce the one-bit slack, but the transfer must match the theorem actually
formalized.

\subsection{Multiplication and magnitude bounds}

Multiplication needs magnitude information.  With
\(|x|,|y|\le M\), \(|x-a|\le\delta\), and \(|y-b|\le\delta\), one has
\[
 |xy-ab|
 \le |x|\,|y-b|+|b|\,|x-a|.
\]
Since \(|b|\le M+\delta\), it suffices to choose \(\delta\) from an explicit
function of \(p\) and a bound code for \(M\).  Accordingly, a multiplication
node transfers output precision to:
\begin{itemize}
  \item higher input precisions;
  \item magnitude-bound resources for both inputs;
  \item any arithmetic overflow or denominator conventions used by the code.
\end{itemize}
A magnitude bound is not silently inferred from existence of the numbers.  It
is a producer theorem or part of the representation contract.

\subsection{Reciprocal and division}

A proof of \(x\ne0\) is insufficient for a finite error modulus of reciprocal.
One needs quantitative apartness, for example
\[
  |x|\ge 2^{-a}.
\]
If \(|x-y|\le 2^{-b}\) and the error is small enough to preserve a lower bound
on \(|y|\), then
\[
  \left|\frac1x-\frac1y\right|
  =\frac{|x-y|}{|x|\,|y|}
\]
has an explicit bound depending on \(a\) and \(b\).  DARC records this
quantitative precondition as a guard coordinate,
\[
  \Guard(\text{denominator},a).
\]
The associated precision transfer has the schematic form
\[
  \Prec(1/x)\ge p
  \Longrightarrow
  \Prec(x)\ge p+2a+c,
  \quad \Guard(|x|\ge2^{-a}).
\]
The constant \(c\) depends on the exact coding and safety slack.  It must be
proved in the local theorem, not guessed by the analysis.

The same principle covers matrix pivots, normalization, logarithms, and
analytic continuation across singular loci.

\subsection{Roots, branches, and collisions}

Ordered root labels are discontinuous at collisions.  A finite approximation
proof should therefore export one of:
\begin{enumerate}
  \item symmetric polynomial observations;
  \item a multiset or cluster code;
  \item a labeled root under an explicit separation guard;
  \item a local branch continued along a finite path with compatibility and
        margin certificates.
\end{enumerate}
A value query for a symmetric observable need not activate branch labels or a
separation guard.  A query that names ``the third root'' must activate both.
This is a canonical example of the benefit of joint certificate/guard demand
analysis.

Spectral theory has the same shape.  At repeated eigenvalues, invariant
subspaces or projectors are stable; individual eigenvectors require a gap,
normalization, and often an orientation convention.  DARC's resource transfer
makes the stronger assumptions appear only for the stronger output.

\subsection{Finite nets instead of abstract compactness}

When an argument uses compactness to select finite data, the constructive
interface is:
\begin{enumerate}
  \item a primitive-recursive finite net at requested scale;
  \item a coverage theorem;
  \item a modulus transferring local estimates from net points to the whole
        domain;
  \item a finite verification of the desired property on the net.
\end{enumerate}
The net and coverage theorem are producer nodes.  If the repository packet has
only an abstract compactness argument, DARC leaves the effective net or
coverage modulus residual instead of presenting the finite check as a complete
proof.

\subsection{Series, integrals, and tail moduli}

A finite partial sum identity is often easy in \(\iPRA\).  An analytic value
requires a tail theorem.  Let \(S_N\) be a primitive-recursive partial sum and
suppose a producer gives
\[
  N\ge M(p,\vec u)
  \Longrightarrow
  |S_N-L|\le 2^{-p}.
\]
Then the requested precision transfers to a cutoff \(M(p,\vec u)\) and the
arithmetic proof verifies the finite sum.  If \(M\) is not proved, the final
endpoint remains conditional on the exact tail formula.  It is incorrect to
replace a convergence assertion by a host-language numerical experiment.

For integration, the same pattern uses a quadrature rule, mesh modulus,
finite error estimate, and endpoint/magnitude guards.  For special functions,
one should separate exact algebraic identities of finite truncations from
analytic continuation and asymptotic remainder producers.

\subsection{Metastability versus convergence rates}

A metastable theorem has finite form
\[
  \forall p,g\;\exists N\le B(p,g)\;
  \forall i,j\in[N,N+g(N)]\;d(x_i,x_j)\le 2^{-p}.
\]
For fixed bounds this can be compiled to a bounded strategy theorem.  It does
not imply a global Cauchy modulus \(M(p)\).  DARC keeps the finite-window
endpoint unless a separate tail argument produces \(M\).  The demand system
must not add a fictitious rate coordinate merely because repeated finite
queries are available.

\subsection{Diophantine finiteness and height bounds}

A theorem that a set of solutions is finite does not provide a number \(H\)
such that every solution has height below \(H\).  A complete enumeration or
least-solution search requires such a cutoff.  DARC splits the task into:
\begin{enumerate}
  \item a primitive-recursive verifier for a proposed solution;
  \item finite enumeration below a supplied height;
  \item a producer theorem asserting completeness of that height.
\end{enumerate}
The first two can be fully compiled.  The third remains residual for general
Faltings-, Siegel-, or Mordell--Weil-type applications unless an effective
bound is actually available.

\subsection{Quantitative transfer as an auditable interface}

Each quantitative node should publish:
\begin{enumerate}
  \item the representation and error predicate;
  \item a monotone resource transfer;
  \item an ordinary local derivation proving that the transferred resources
        suffice;
  \item the conventions at singular or failed cases;
  \item a status for each producer theorem.
\end{enumerate}
This interface is analogous to a function type with preconditions and a
modulus.  It is the main mechanism by which DARC turns informal analytic
reasoning into a sequence of finite proof obligations.
